% Chapter 2, section 2: Create and configure queues
% Write in plain LaTeX. Code of a guide step: \guide{<step>}; commands: \begin{shell}...\end{shell};
% a screenshot: \screenshot{file}{caption}; a gap in the draft: \placeholder{...}.
% See book/latex/queen.sty for the full list.

\section{Create and configure queues}

A queue is a named set of partitions that share one configuration. You never have to create one:
the first push that names a queue creates it, with the defaults. Configure it when you want other
settings.

\queenfig{queue-anatomy}{A queue: options shared by its partitions, each an ordered list of messages.}

\guide{configure}

\why{a queue created by a push is good enough to start, and most queues never need more. When they
do, the settings live on the queue, in one place, instead of in every worker.}

\subsection{Every option}

Here is a queue that sets most of them, and the table below says what each one does.

\guide{configure-jobs}

\begin{options}{Option}{Default}{What it does}
leaseTime & 60 s (300 s from \texttt{create()}) & How long a worker holds the messages it took before they go to another worker. \\
retryLimit & 3 & How many failed attempts a message gets. After the last one it goes to the dead-letter queue. \\
retryDelay & 1000 ms & How long a failed message waits before it comes back. Its partition waits with it, for that consumer group. \\
deadLetterQueue, dlqAfterMaxRetries & true & File a message that ran out of retries in the dead-letter queue. With both off, it is dropped. \\
delayedProcessing & 0 s & How long every new message waits after its push before anyone can take it. \\
windowBuffer & 0 s & How long new messages wait in their partition, so a burst is handed over as one batch. \\
dedupWindowSeconds & 3600 s & How long a partition remembers \texttt{transactionId}s. \texttt{0} turns deduplication off. \\
retentionEnabled & false & Must be true for the two options below to do anything. \\
retentionSeconds & 0 & Remove messages older than this, consumed or not. \\
completedRetentionSeconds & 0 & Remove messages older than this that every consumer group has already acked. \\
maxWaitTimeSeconds & 0 & Remove messages older than this, consumed or not. Works without \texttt{retentionEnabled}. \\
encryptionEnabled & false & Encrypt payloads on disk. Every node needs \texttt{QUEEN\_ENCRYPTION\_KEY}. \\
namespace, task & from the name & Labels that group queues in the dashboard and in a discovery pop. \\
\end{options}

\paragraph{Leases.} A worker that takes messages holds them under a lease. Until it acks them or the
lease runs out, no other worker of the same consumer group gets that partition.

\why{a worker can crash in the middle of its work. The lease is how its messages get back to
someone else without a human noticing. Set it longer than your slowest handler, or renew it while
the handler runs (see Pop messages).}

\paragraph{Retries.} A failed message comes back \texttt{retryLimit} times, each time after
\texttt{retryDelay}, and then goes to the dead-letter queue.

\why{most failures are temporary: a database that restarts, a provider that times out. Trying again
later fixes them. Waiting between tries gives the other side time to recover, and a message that
can never succeed stops after a few tries instead of blocking its partition forever.}

\paragraph{Delays.} \texttt{delayedProcessing} holds every new message of the queue for a while.
\texttt{windowBuffer} holds a group's new messages until the burst is over, so the handler gets them
together.

\why{some work should not start at once (a reminder sent five minutes after a signup), and some
work is cheaper in bulk (one database write for a burst of fifty updates). For a delay that is
different for each message, use a timer instead (see KV, Locks and Timers).}

\paragraph{Deduplication.} Inside the dedup window, a second push with the same
\texttt{transactionId} in the same partition is not written.

\why{producers retry when they do not get an answer, and a retry must not create a second order.}

\paragraph{Retention.} By default Queen keeps every message, and the disk grows. Retention runs only
when \texttt{retentionEnabled} is true; then \texttt{retentionSeconds} removes old messages, read or
not, and \texttt{completedRetentionSeconds} removes old messages that every group has finished.
\texttt{maxWaitTimeSeconds} removes old messages without needing \texttt{retentionEnabled}. A node
whose disk is 85\% full refuses the writes that would grow it.

\why{an acked message is not deleted: it stays in its partition so that a new consumer group, an
audit or a fixed bug can read it again. Retention is how you decide how long that history is worth
its disk.}

\paragraph{Encryption.} With \texttt{encryptionEnabled}, payloads are encrypted on disk.

\why{some payloads hold personal data that must not sit on a disk in clear text.}

\subsection{Three traps}

\begin{itemize}
\item \texttt{config(...).create()} in the JavaScript client sends your options together with the
client's own defaults (a 300-second lease, encryption off, no retention). Calling it on an existing
queue resets every option you leave out. Name every option you care about.
\item On a node without \texttt{QUEEN\_ENCRYPTION\_KEY}, an encrypted queue stores plain text, and
nothing warns you.
\item \texttt{priority}, \texttt{ttl}, \texttt{maxSize} and \texttt{minPopWaitTime} are accepted and
stored, and do nothing.
\end{itemize}

On the site, with the code in all six SDK languages: \href{https://queenmq.com/writing-applications/queues/}{queenmq.com/writing-applications/queues}.
